\color{unidarkgreen}{\subsection{Токовая защита нулевой последовательности}}
\color{black}

\begin{enumerate}

\item
Для ШР предусматривается два комплекта токовой защиты нулевой последовательности, предназначенных для резервирования действия основных быстродействующих защит при внутренних земляных видах коротких замыканий. Ступени ТЗНП обоих комплектов имеют идентичную логику работы, ниже приводится описание работы ступени ТЗНП-1 ВВ.

\item
Вывод ступени производится накладкой \textbf{«ТЗНП-1 ВВ»}.

\item
Измерительные цепи первого комплекта ТЗНП включены на токи фаз А, В, С со стороны высоковольтных вводов ШР. Защита реагирует расчетное значение утроенного тока нулевой последовательности, определяемого по формуле:

\begin{equation}
\vec{3I_{0}} = \vec{I_{A}} + a^2 \cdot \vec{I_{B}} + a \cdot \vec{I_{C}}
\end{equation}

\item
Второй комплект ТЗНП может быть включен на токи фаз А, В, С со стороны нейтральных выводов ШР или к измерительному трансформатору току, установленному в цепи заземления нейтрали ШР. Данный комплект обладает высокой чувствительностью к однофазным коротким замыканиям вблизи нейтрали.

\item
Каждая из ступеней защиты может выполнена с блокировкой по второй гармонике (отношению уровня второй гармоники к первой). Данная блокировка реализована на блоке «ОВ БНТ ВВ». Выводится данная блокировка накладкой \textbf{«Режим БНТ ТЗНП-1 ВВ»}.

\item
Защита имеет регулируемые значения по току и времени срабатывания, задаваемые параметрами \textbf{«3I0ср ТЗНП-1 ВВ»} и \textbf{«Тср ТЗНП-1ВВ»} соответственно.

\item 
Уставки ТЗНП ВВ приведены в таблице \ref{app_set:tznpvv}, уставки ТЗНП НВ приведены в таблице \ref{app_set:tznpnv}. Алгоритм работы ТЗНП-1 ВВ приведен на рисунке \ref{pict_tznp:TZNP}, алгоритмы работы ТЗНП-2 ВВ, ТЗНП-1 НВ, ТЗНП-2 НВ выполнены аналогично.

\medskip
\begin{figure}[H]
\centering
\includegraphics[width=1\textwidth,height=1\textheight,keepaspectratio]{\fbpath/11.01 ШР/171.0241 ТЗНП/img/TZNP.pdf}
\captionof{figure}{Функциональная схема логики ТЗНП-1 ВВ}\label{pict_tznp:TZNP}
\end{figure}


\end{enumerate}